% ============================================================
% Preamble chung cho cả hai bản báo cáo — biên dịch bằng XeLaTeX
% Bố cục một cột, giữ font / lề / kiểu tiêu đề theo phong cách ACL
% ============================================================
\usepackage{amsmath,amssymb}
\usepackage{fontspec}
\setmainfont{TeX Gyre Termes}      % giống Times (font của ACL), đủ dấu tiếng Việt
\setsansfont{TeX Gyre Heros}
\setmonofont{TeX Gyre Cursor}
\usepackage{babel}
\babelprovide[import,main]{vietnamese}   % cách khuyến nghị cho XeLaTeX

\usepackage[a4paper,margin=2.5cm]{geometry}
\usepackage{graphicx}
\graphicspath{{figures/}}
\usepackage{booktabs,multirow,array,tabularx,makecell}
\usepackage[table]{xcolor}
\usepackage{caption}
\captionsetup{font=small,labelfont=bf,labelsep=period}
\usepackage{enumitem}
\setlist{itemsep=2pt,topsep=4pt}
\usepackage{titlesec}
\titleformat{\section}{\large\bfseries}{\thesection}{0.6em}{}
\titleformat{\subsection}{\normalsize\bfseries}{\thesubsection}{0.6em}{}
\titleformat{\subsubsection}[runin]{\normalsize\bfseries}{\thesubsubsection}{0.5em}{}[.]
\titlespacing*{\section}{0pt}{14pt}{6pt}
\titlespacing*{\subsection}{0pt}{10pt}{4pt}
\usepackage{tikz}
\usetikzlibrary{positioning,arrows.meta,fit}
\usepackage[round,authoryear]{natbib}
\usepackage[colorlinks=true,linkcolor=blue!50!black,citecolor=blue!50!black,urlcolor=blue!50!black]{hyperref}

\setlength{\parskip}{3pt}
\renewcommand{\arraystretch}{1.15}

% Tên chú thích tiếng Việt (đặt lại cho chắc, phòng khi babel không đổi)
\addto\captionsvietnamese{%
  \renewcommand{\figurename}{Hình}%
  \renewcommand{\tablename}{Bảng}%
  \renewcommand{\contentsname}{Mục lục}%
  \renewcommand{\refname}{Tài liệu tham khảo}%
  \renewcommand{\abstractname}{Tóm tắt}%
  \renewcommand{\appendixname}{Phụ lục}%
}

% Chèn hình; nếu chưa upload file .png lên Overleaf thì hiện khung báo thiếu thay vì lỗi biên dịch
\newcommand{\hinh}[2][\linewidth]{%
  \IfFileExists{figures/#2}{\includegraphics[width=#1]{#2}}{%
    \fbox{\parbox[c][3cm][c]{0.8\linewidth}{\centering\small Thiếu file \texttt{figures/#2}}}}}

\newcommand{\RQ}[1]{\textbf{RQ#1}}
\newcommand{\metamodel}{meta-model}
